\section{O Método Clássico e o Teorema Central do Limite}

\begin{frame}{O Teorema Central do Limite (TCL)}
    \begin{block}{Definição}
        Independentemente da distribuição da população original, a distribuição da média amostral $\bar{X}$ se aproxima de uma distribuição normal à medida que o tamanho da amostra $N$ cresce.
    \end{block}
    
    \begin{block}{Intervalo de Confiança (IC) Clássico}
        O intervalo de 95\% de confiança para a média é dado por:
        $$IC = \bar{X} \pm 1.96 \times \frac{s}{\sqrt{N}}$$
        Onde $s$ é o desvio padrão amostral e $\frac{s}{\sqrt{N}}$ é o erro padrão (SE).
    \end{block}
\end{frame}

\begin{frame}{Limitações da Abordagem Clássica}
    \begin{alertblock}{Por que o TCL nem sempre é suficiente?}
        \begin{itemize}
            \item \textbf{Amostras pequenas:} Se a população for muito assimétrica e $N$ for pequeno, a aproximação normal falha.
            \item \textbf{Estatísticas complexas:} E se quisermos construir um intervalo de confiança para a \textbf{mediana}, \textbf{variância} ou uma razão entre variáveis? A matemática clássica torna-se extremamente complexa ou indisponível.
            \item \textbf{Outliers:} Presença de valores extremos distorce o desvio padrão e o erro padrão.
        \end{itemize}
    \end{alertblock}
    \centering
    \textbf{Solução em Ciência de Dados: Métodos Computacionais (Bootstrap)}
\end{frame}

\begin{frame}{Exemplo Matemático: IC via TCL}
    \begin{block}{Problema}
        Amostra de $N=36$ tempos de atendimento com média $\bar{X}=100\text{ s}$ e desvio padrão $s=15\text{ s}$. Queremos o IC de 95\% para a média populacional $\mu$.
    \end{block}
    \begin{block}{Cálculo Analítico (TCL)}
        \begin{itemize}
            \item \textbf{Erro padrão da média ($SE$):} $SE = \frac{s}{\sqrt{N}} = \frac{15}{\sqrt{36}} = 2.5\text{ s}$
            \item \textbf{Margem de erro (para $z = 1.96$):} $1.96 \times 2.5 = 4.9\text{ s}$
            \item \textbf{Intervalo de Confiança (95\%):} $IC = 100 \pm 4.9 = [95.1\text{ s}, 104.9\text{ s}]$
        \end{itemize}
    \end{block}
\end{frame}

\begin{frame}[fragile]{Demonstração em Python: IC Clássico}
    \begin{lstlisting}[language=Python]
import numpy as np
import scipy.stats as ss

# Dados amostrais simulados (N=36, media=100, std=15)
np.random.seed(42)
dados = np.random.normal(loc=100, scale=15, size=36)

# Calculos intermediarios
n = len(dados)
media = np.mean(dados)
desvio = np.std(dados, ddof=1) # Desvio padrao amostral
erro_padrao = desvio / np.sqrt(n)

# Encontrando o z critico para 95%
z_critico = ss.norm.ppf(0.975)

# Limites do intervalo
lim_inf = media - z_critico * erro_padrao
lim_sup = media + z_critico * erro_padrao
print(f"IC Classico: [{lim_inf:.2f}, {lim_sup:.2f}]")
# Saida aproximada: IC Classico: [94.59, 104.07]
    \end{lstlisting}
\end{frame}

